\documentclass[a4paper,11pt]{article}
\usepackage[T1]{fontenc}
\usepackage[utf8]{inputenc}
\usepackage{t1enc}
\usepackage[dvips]{color}
\usepackage{amsmath,amssymb,makeidx,stmaryrd}
\usepackage[normalem]{ulem}
\usepackage{fancybox}
\definecolor{gris}{gray}{0.9}
\usepackage{pst-plot,pst-tree,pstricks,pst-node,pst-text}
\usepackage{pst-eucl}
\usepackage{pstricks-add}
\usepackage[left=3.5cm, right=3.5cm, top=3cm, bottom=3cm]{geometry}
\newcommand{\vect}[1]{\overrightarrow{\,\mathstrut#1\,}}
\def\bas{\fcolorbox{gris}{gris}{\parbox{0.986\textwidth}{
	\color{red}\footnotesize 1991\hfill Page : \thepage/1}}}
\def\haut{\fcolorbox{gris}{gris}{\parbox{0.986\textwidth}{
	\color{red}
	\bf Concours Général de Mathématiques \hfill 1991
	}}}
\makeatletter
\def\ps@cg{%
\let\@oddhead\haut%
\let\@evenhead\haut%
\let\@oddfoot\bas
\let\@evenfoot\bas}
\makeatother
\newcommand{\R}{\mathbb{R}}
\newcommand{\N}{\mathbb{N}}
\newcommand{\D}{\mathbb{D}}
\newcommand{\Z}{\mathbb{Z}}
\newcommand{\Q}{\mathbb{Q}}
\newcommand{\C}{\mathbb{C}}
\pagestyle{cg}
\renewcommand{\theenumi}{\textbf{\arabic{enumi}}}
\renewcommand{\labelenumi}{\textbf{\theenumi.}}
\renewcommand{\theenumii}{\textbf{\alph{enumii}}}
\renewcommand{\labelenumii}{\textbf{\theenumii.}}
\parindent0pt
%\everymath{\displaystyle}
\usepackage[french]{babel}

\DecimalMathComma
\usepackage[np]{numprint}
\begin{document}

\begin{center}\textbf{{\color{blue} Exercice I}}
\end{center}

\begin{enumerate}
\item Soit $\left(x_n\right)_{n\in \N}$ une suite de nombres réels telle que, pour tout nombre entier naturel $n$:

\[x_0 ^3 + x_1 ^3 +...+ x_n^3 = \left(x_0 + x_1 + \ldots + x_n\right)^2.\]

Montrer que, pour tout entier naturel $n$, il existe un entier naturel $m$ tel que :

\[x_0 + x_1 + \ldots + x_n = \dfrac{m(m + 1)}{2}.\]

\item Si $n$ et $p$ sont deux nombres entiers naturels non nuls, on pose :

\[S_{n, p} = 1^p + 2^p + \ldots + n^p.\]

Déterminer les entiers naturels non nuls $p$ tels que, quel que soit l'entier naturel non nul $n$, \:$S_{n, p}$ soit le carré d'un nombre entier naturel.
\end{enumerate}

\begin{center}\textbf{{\color{blue} Exercice II}}
\end{center}

À tout nombre entier naturel non nul $n$, on associe l'application $f_n$ de la variable réelle $x$, définie pour $x \geqslant n$ par :

\[f_n(x) = \sqrt{x-n} + \sqrt{x-n+1} + \ldots +\sqrt{x-1} + \sqrt{x} +\sqrt{x+1} +\ldots  + \sqrt{x+n} - (2n+1)\sqrt{x}\]

\begin{enumerate}
\item Dans cette question, l'entier $n$ est fixé.

Montrer que $f_n$ est croissante et que $\displaystyle\lim_{x \to + \infty}f_n(x) = 0$.
\item Déterminer la limite de la suite de terme général $f_n(n)$.
\end{enumerate}

\begin{center}\textbf{{\color{blue} Exercice III}}
\end{center}

Soit S un point fixe d'une sphère fixe $(\sum)$ de centre $\Omega$.

On considère les tétraèdres SABC inscrits dans la sphère $(\sum)$ et dont les arêtes issues de S sont deux à deux orthogonales.

\medskip

\begin{enumerate}
\item Montrer que les plans (ABC) passent par un point fixe.
\item Pour un tel tétraèdre SABC, le point S et le centre $\Omega$ de la sphère $(\sum)$ se projettent orthogonalement sur le plan (ABC) respectivement en H et O. 

On note $R$ le rayon du cercle circonscrit au triangle ABC.

Démontrer que : $R^2 = \text{OH}^2 +2\text{SH}^2$.
\end{enumerate}

\begin{center}\textbf{{\color{blue} Exercice IV}}
\end{center}

Soit $p$ un nombre entier naturel et $n = 2^p$.

On considère les parties $A$ de l'ensemble $E = \{1, 2, \ldots , n\}$ possédant la propriété suivante :

\begin{center}
si $x$ appartient à $A$, alors $2x$ n'appartient pas à $A$.\end{center}

Déterminer le nombre maximal d'éléments d'un telle partie A.

\begin{center}\textbf{{\color{blue} Exercice V}}
\end{center}

\begin{enumerate}
\item On considère l'application $P$ de $\C$ dans $\C$ définie par :

\[P(z) = z^5 + a_4z^4 + a_3z^3 + a_2z^2 + a_1z\]

où $a_1, a_2, a_3$ et $a_4$ sont quatre nombres complexes donnés.

On pose $w_l = \text{e}^{2\text{i}l\pi/5}$ où $l$ désigne un nombre entier compris entre 0 et 4.

Montrer que $P\left(w_0\right) + P\left(w_1\right) + P\left(w_2\right) + P\left(w_3\right) + P\left(w_4\right) = 5$.
\item Soient A$_1$, A$_2$ A$_3$, A$_4$, A$_5$ cinq points du plan. 

On construit un pentagone régulier inscrit dans le cercle de centre A$_1$ et de rayon donné $R$.

Démontrer qu'il existe un sommet S du pentagone tel que:

\[\text{SA}_1 \cdot \text{SA}_2 \cdot\text{SA}_3 \cdot\text{SA}_4 \cdot\text{SA}_5 \geqslant R^5.\]
\end{enumerate}
\end{document}